\subsubsection{Robustness}

\begin{theorem} \label{thm:perfect-robust}
  The generic construction $\dkg$ is weakly robust in the honest majority setting
  ($t-1 < n/2$) against
  an adversary $\adv$ controlling up to $t-1$ participants in the random oracle model,
  given a zero-indexed and secure \AVSS $\avss$
  and a binding \ENIKE $\enike$.

  Concretely, for every adversary $\adv$ that attacks the scheme,
  there exist adversaries $(\bdv_1, \bdv_2)$ that run in about the same time as $\adv$ such that
  \begin{align} \label{eq:wk-rbt-thm}
       \advantwkrobust{\dkg,\adv}(\lambda) \leq & \ \advant{bind}{\enike,\bdv_1}(\lambda) +
       \advantuniqueness{\avss[\Hash_1],\bdv_2}(\lambda).
    \end{align}
\end{theorem}

\begin{proof}
%\setcounter{case}{0}

We prove Theorem~\ref{thm:perfect-robust} by a sequence of games.

\medskip

\noindent \textit{\underline{Game 0.}}
This is the DKG weak robustness game as in Figure~\ref{fg:robustness},
applied to the construction in Figure~\ref{fg:generic-perfect}.

Let $W_0$ be the event that $\adv$ wins in Game~0. Then,

\[
    \advant{wk-rbst}{\dkg,\adv}(\lambda) = \Pr[W_0]
\]


%Let $\cdv$ be the algorithm that simulates the robustness game to $\adv$.
%We now describe how $\cdv$ simulates the game.
%
%\medskip
%
%\noindent \textit{Setup.}
%To begin, $\cdv$ initializes $Q_1 \gets \emptyset$ to simulate $\Hash_1$
%and $Q_2 \gets \emptyset$ to simulate $\Hash_2$,
%in addition to initializing each honest party as shown in Figure~\ref{fg:robustness}.
%$\cdv$ handles $\adv$'s random oracle queries by lazy sampling, as follow:
%
%  \begin{itemize}
%    \item[$\underline{\Hash_1}$]: When the adversary queries $\Hash_1$ on inputs $(\pubaggregateset_i, \auxstring_i)$,
%      the environment checks if it is in $Q_1$, and if so, returns the corresponding $\tweak_i$.
%      Otherwise, the environment randomly samples $\tweak_i$ from its respective domain,
%      sets $Q_1[(\pubaggregateset_i, \auxstring_i)] = \tweak_i$, and then returns $\tweak_i$.
%    \item[$\underline{\Hash_2}$]: When the adversary queries $\Hash_2$ on inputs $\{\blinding_1,\ldots,\blinding_n \}$,
%      the environment checks if it is in $Q_2$, and if so, returns the corresponding $\auxstring_i$.
%      Otherwise, the environment randomly samples $\auxstring_i$ from its respective domain,
%      sets $Q_2[\{\blinding_1,\ldots,\blinding_n \}] = \auxstring_i$, and then returns $\auxstring_i$.
%  \end{itemize}
%
%$\cdv$ performs all signing oracle queries honestly.

\medskip

\noindent \textit{\underline{Game 1.}}
The only difference in Game~1 is if $\adv$ outputs a tuple $(\commitment_j, \pk_j,\vsscommit_j,\witness_j), j \in \corrupt$
where $(\commitment_j, \pk_j,\vsscommit_j)$ are output in $\dkgkeygen_0$
and $\witness_j$ is output in $\dkgkeygen_2$
such that Equation~\ref{eq:game1-robust-condition} holds,
then the environment aborts.

\begin{align} \label{eq:game1-robust-condition}
  \begin{split}
    &\ \avss[\Hash_1].\derivepshare(0, \vsscommit_j) = \pk, \text{ and } \\
    &\ \enike.\verify(\witness_j, \commitment_j) \rightarrow 1, \text{ but } \\
    &\ \enike.\getsharedkey(\witness_j, \pk_j) \neq \enike.\getsharedkey(\secretone_j, \commitment_j)
  \end{split}
\end{align}

However, if the \ENIKE is binding and Assumption~\ref{ass:avss-enike} holds,
then Game~1 and Game~0 are indistinguishable to $\adv$.

\medskip

\noindent \textit{Reduction to \ENIKE Binding.}
Let $\bdv_1$ be an adversary playing against the \ENIKE binding game.
%\douglas{The general structure of this proof doesn't seem right at all.  The \ENIKE binding game isn't a distinguishing game, there's no hidden bit $b$.  The \ENIKE binding game is a win/lose game.  Now you can certainly still do a reduction to the binding game, and just observe that when the abort event happens in Game 1 corresponds exactly to when the reduction wins the \ENIKE binding game.}
% CK- changed it to this, hope it looks right now
$\bdv_1$ simulates Game~1 to $\adv$ in exactly the same way as Game~0.
However, in $\finalize$,
if $\adv$ outputs some $\witness_j$ with respect to the commitments $(\pk_j, \commitment_j, \vsscommit_j)$ output in $\dkgkeygen_0$ such that Equation~\ref{eq:game1-robust-condition} holds,
then $\bdv_1$ recovers $\secretone_j$,
and outputs $(\secretone_j, \pk_j, \witness_j, \commitment_j)$ as its output to the \ENIKE binding game.

$\bdv_1$ recovers $\secretone_j$ by performing $\secretone_j \gets \avss[\Hash_1].\recover(t, \recoveryset_j)$,
where
$\recoveryset_j = \{(k,\share_{jk}) \}_{k \in \honest}$.
$\bdv_1$ can recover $\secret_j$ because the checks in $\dkgkeygen_1$ completed successfully
(else the honest participants would have aborted the protocol in $\dkgkeygen_1$ and $\dkgkeygen_2$),
then each share in $\recoveryset_j$ is a valid share of the secret $\secretone_j$ committed to by $\vsscommit_j$ and $\pk_j$.
Further, because $\bdv_1$ simulates at least $t$ honest players,
it has a sufficient number of shares to perform this recovery step.

Hence, when $\adv$ distinguishes Game~1 from Game~0 (i.e., by causing Game~1 to abort),
then $\bdv_1$ wins.

\medskip

\noindent \textit{Difference between Game~0 and Game~1.}
Let $W_1$ be the event that $\adv$ wins in Game~1. Then,
\begin{equation} \label{eq:game1-robust}
  \bigabs{ \Pr[W_1] - \Pr[W_0]  } \leq \advant{bind}{\enike,\bdv_1}(\lambda)
\end{equation}

\medskip


\noindent \textit{\underline{Game 2.}}
The only difference in Game~2 is that in $\finalize$,
if Equation~\ref{eq:game2-robust-condition} holds,
then the environment aborts.
\begin{align} \label{eq:game2-robust-condition}
  \begin{split}
    &\ \sk \gets \avss.\recover(t, \{(i, \sk_i) \}_{i \in \honest}), \text{ but} \\
    &\ \enike.\verify(\sk, \pk) \rightarrow 0
  \end{split}
\end{align}

However, if the \AVSS is unique, then Game~1 and Game~2 are indistinguishable.

\medskip

\noindent \textit{Reduction to \AVSS Uniqueness.}
Let $\bdv_2$ be an adversary playing against the \AVSS uniqueness game.

$\bdv_2$ follows the protocol as defined in Figure~\ref{fg:generic-perfect} honestly for all parties $i \in \honest$.
However, in $\finalize$, $\bdv_2$ recovers $\sk$ and tests to see if Equation~\ref{eq:game2-robust-condition} holds.
$\bdv_2$ can recover $\sk$ without the involvement of $\adv$ because it simulates at least $t$ honest parties.

Recall that

\[
\sk_i \gets \avss[\Hash_1].\aggregatepriv(i, \{\share_{i,j} \}_{j \in \set{n}}, \{ \vsscommit_j \}_{j \in \set{n}}, \auxstring)
\]

and

\[
\avsscommit \gets \avss[\Hash_1].\aggregatepub(\{ \vsscommit_j \}_{j \in \set{n}}, \auxstring)
\]

The only values that differ between each honest party when performing this aggregation step are the sets $\{\share_{i,j} \}_{j \in \set{n}}$.

If Equation~\ref{eq:game2-robust-condition} holds,
then this means that at least one honest party $k \in \honest$ received a share $\share_{j,k}$ from a corrupted party $j \in \corrupt$
such that $\avss[\Hash_1].\verify(k, \share_{j,k}, \vsscommit_j) = 1$
but where Equation~\ref{eq:g2-robust-recovery} holds.
\begin{align} \label{eq:g2-robust-recovery}
  \begin{split}
 & \avss[\Hash_1].\recover(t, \{(\ell,\share_{j,\ell}) \}_{\ell \in \honest}) \neq \\
 & \avss[\Hash_1].\recover(t, \{(\ell,\share_{j,\ell}') \}_{\ell \in \honest})
  \end{split}
\end{align}
%\ian{API for $\avss.\recover$ above.  But worse, doesn't the second recovery set only have $t-1$ elements? $\recover$ isn't even defined for such an input?}
%\chelsea{The adversary trivially cannot win if it chooses $\lvert \honest \rvert = t$}


where $\share_{j, \ell} = \share_{j, \ell}'$,
except for $\ell = k$.
Let $\recoveryset_1 = \{(\ell,\share_{j,\ell}) \}_{\ell \in \honest}$ and
$\recoveryset_2 = \{(\ell,\share_{j,\ell}') \}_{\ell \in \honest}$.
When this event has occurred,
$\bdv_2$ submits $(t, \recoveryset_1, \recoveryset_2, \vsscommit_j)$ as its output to the \AVSS uniqueness game.
%\ian{The $t$ and $\ell$ should both be omitted, right? (And check your paren balencing.) And this is the ``uniqueness'' game, not the ``security'' game, right?}
%\chelsea{The adversary needs to return t, two recovery sets (which are the identifier/share pairs) and the commitment}
When $\adv$ can distinguish Game~2 from Game~1,
then $\bdv_2$ wins the \AVSS uniqueness game.

\medskip

\noindent \textit{Difference between Game~1 and Game~2.}
Let $W_2$ be the event that $\adv$ wins in Game~2. Then,
\begin{equation} \label{eq:game2-robust}
  \bigabs{ \Pr[W_2] -  \Pr[W_1]  } \leq \advant{unq}{\avss[\Hash_1],\bdv_2}(\lambda)
\end{equation}



\medskip

\noindent \textit{Finishing the Proof.}
Recall that in the weak robustness game,
the adversary has four opportunities to win,
assuming all honest parties have not aborted the protocol.
First, if any honest party has a disjoint view of $\pk$,
second, if any honest party has a disjoint view of $\qual$,
and third, if any honest party has a disjoint view of $\auxinfo$.
Finally, the adversary wins if the $\sk$ that is recovered from the honest parties' shares is invalid.

Because we assume a broadcast channel,
if any honest party aborts in either $\dkgkeygen_1$ or $\dkgkeygen_2$,
then all honest parties will abort in $\dkgkeygen_2$.
And so if any honest party proceeds to $\finalize$,
then all honest parties will proceed to $\finalize$,
and all honest parties will set their status to $\accept$.
While honest parties might abort (meaning that strong robustness cannot be fulfilled),
honest parties will consistently either end in an aborted or accepting status,
meaning that the weak robustness game can be fulfilled.
And so the only opportunity for the adversary to win against the generic
construction in the weak robustness game is to disrupt consistency of $\pk,\qual,\auxinfo$,
or to make $\dkg[\Hash].\recover$ output $\fail$.

Again, because we assume a broadcast channel,
all parties will maintain the same view of $\{(\witness_j,\commitment_j,\pk_j,\vsscommit_j) \}$.
And so all parties will determine the set $\qual$ in $\finalize$ using the same information.
Hence, the adversary has no advantage in forcing a disjoint $\qual$.

From Game~1,
we know that all honest parties derive the same view of $\{ \blinding_j \}$ such that $\enike.\getsharedkey(\witness_j,\pk_j) \rightarrow \blinding_j$.
Because the tweak value $\tweak$ is derived from $\Hash_2$ using the values $\{\blinding_j \}_{j \in \set{n}}$,
then all parties will obtain the same $\tweak$.
Because $\PK$ is derived using $\tweak$ and the set $\{\vsscommit_j \}_{j \in \set{n}}$,
then the adversary has no advantage in forcing a disjoint $\PK$.
Similarly,
the adversary has no advantage in forcing a disjoint $\auxinfo = \{\PK_i \}_{i
\in \qual}$, which are also derived from $\tweak$ and the set $\{\vsscommit_j \}_{j \in \set{n}}$.

Finally, from Game~2,
we know that $\dkg[\Hash].\recover$ will output a $\sk$ that is valid with respect to $\pk$.

The adversary in Game~2 then has no advantage in winning the DKG weak robustness game.
Concretely, combining the advantages in
Equations~\ref{eq:game1-robust} and \ref{eq:game2-robust} gives Equation~\ref{eq:wk-rbt-thm}.
This completes the proof. \qed
\end{proof}

